% ============================================================
% Capitolo 3 — Una sessione di lavoro reale
% ============================================================

\chapter{Una sessione di lavoro reale}
\label{ch:sessione-reale}

È martedì mattina. Lorenzo apre il terminale, avvia Claude Code e inizia la sessione su TaskFlow API. Seguiamo ogni passaggio, dal bootstrap al commit finale, per vedere ADLC in azione su un caso concreto.

\begin{notebox}
TaskFlow API ha attraversato le fasi 0--2 nelle settimane precedenti. I requisiti sono stati analizzati, l'architettura definita (schema DB, contratto API con 12 endpoint), le decisioni aperte chiuse. Lorenzo ha aggiornato \texttt{\_CONTEXT.md} a \texttt{Phase: 3-Implementation} al termine del design. Da questo punto inizia il lavoro che seguiamo ora.
\end{notebox}

Il task di oggi: implementare l'endpoint \texttt{POST /tasks} --- il cuore dell'API, quello che crea un nuovo task per un utente autenticato.

\section{Inizio sessione --- il bootstrap}
\label{sec:03-bootstrap}

Lorenzo non spiega nulla all'agente. Chiede solo:

\begin{quote}
``Carica \texttt{\_CONTEXT.md} e confermami il contesto.''
\end{quote}

L'agente risponde:

\begin{lstlisting}[language=,caption={Conferma di bootstrap}]
Context: TaskFlow API | Phase: 3-Implementation |
Task: T-007 POST /tasks endpoint |
Stack: Node.js 22/Fastify 5/PostgreSQL 16 |
Constraints: SEC-01,SEC-02,SEC-03,PERF-01,PERF-02. Proceed?
\end{lstlisting}

In tre secondi, l'agente sa tutto. Lorenzo scrive: \textit{``Procedi.''}

\begin{tipbox}
\textbf{Pattern da seguire:} all'inizio di ogni sessione, chiedi sempre la conferma del contesto prima di qualsiasi istruzione. È il modo per verificare che l'agente abbia letto \texttt{\_CONTEXT.md} correttamente.
\end{tipbox}

\section{Task LOW --- una modifica senza freni}
\label{sec:03-task-low}

Lorenzo comincia con un task leggero: aggiungere JSDoc all'interfaccia \texttt{Task} già esistente. È una modifica LOW-risk: documentazione, nessun impatto sulla logica.

\begin{quote}
``Aggiungi JSDoc all'interfaccia \texttt{Task} in \texttt{src/types/task.ts}. Segui lo stile già usato in \texttt{src/types/user.ts}.''
\end{quote}

L'agente produce il JSDoc coerente con lo stile del progetto e notifica:

\begin{lstlisting}[language=,caption={Notifica per task LOW}]
[LOW] Added JSDoc to Task interface (4 fields documented).
No risk items. Changes in: src/types/task.ts
\end{lstlisting}

Per i task LOW, l'agente esegue e notifica senza aspettare approvazione.

\section{Task MEDIUM --- proposta e approvazione}
\label{sec:03-task-medium}

Lorenzo chiede di implementare l'endpoint principale:

\begin{quote}
``Implementa \texttt{POST /tasks}. L'endpoint riceve titolo, descrizione e scadenza, crea il task per l'utente autenticato, e risponde 201 con il task creato.''
\end{quote}

L'agente classifica come MEDIUM (nuova feature) e risponde con un piano prima di scrivere codice:

\begin{lstlisting}[language=,caption={Piano proposto per task MEDIUM}]
[MEDIUM] Piano per POST /tasks:

1. Route handler in src/routes/tasks.ts
   - Validazione: Zod schema (SEC-01)
   - userId dal JWT, non dall'input (SEC-02)
   - Nessun userId nei log (SEC-03)

2. Service in src/services/taskService.ts
   - createTask(userId, data) → Prisma
   - Errori: 400 / 401 / 500 in formato RFC 7807

3. Test in tests/routes/tasks.test.ts
   - 201 dati validi, 400 titolo mancante,
   - 401 no auth, 400 titolo troppo lungo

Stima: ~80 righe codice + ~60 test. Model Level: 3.
Procedo?
\end{lstlisting}

Lorenzo approva. L'agente implementa citando SEC-02 e SEC-03 nel codice generato --- i vincoli sono attivi, non devono essere ripetuti.

\section{Task HIGH --- l'HALT in azione}
\label{sec:03-task-high}

Lorenzo nota che lo schema Prisma non ha un indice su \texttt{userId} nella tabella \texttt{tasks}. Chiede all'agente di aggiungerlo.

L'agente si ferma:

\begin{lstlisting}[language=,caption={HALT per modifica schema}]
[!] HALT — trigger: schema (HIGH)

File: prisma/schema.prisma
Pattern: **/prisma/schema.prisma

Proposta: aggiunta @@index([userId]) in model Task.
Richiede migration:
  prisma migrate dev --name add-task-userid-index

Rischio: MEDIUM — migration non-destructive.
In produzione: lock temporaneo sulla tabella.

Minimum model level: 5. Conferma esplicita richiesta.
Procedo?
\end{lstlisting}

Lorenzo risponde: \textit{``Confermato. Procedi con schema e migration.''}

L'agente aggiunge l'indice e genera il file di migration. Il processo di HALT ha reso espliciti i dettagli che meritavano attenzione.

\section{Il checkpoint}
\label{sec:03-checkpoint}

Dopo circa un'ora di lavoro, Lorenzo manda un checkpoint:

\begin{quote}
\texttt{@checkpoint}
\end{quote}

L'agente risponde:

\begin{lstlisting}[language=,caption={Output di @checkpoint}]
Checkpoint — TaskFlow API | 2026-06-27 10:47

Completato:
- POST /tasks: route, service, Zod, RFC 7807, 4 test
- Migration: add-task-userid-index

In corso: niente di sospeso. Blocchi: nessuno.

Aggiornamento _CONTEXT.md consigliato:
  Active Task → T-007.2 GET /tasks
  Active Task Token Est. → 5000/1500/6500

Applico gli aggiornamenti?
\end{lstlisting}

Lorenzo dice sì. L'agente produce il blocco da incollare in \texttt{\_CONTEXT.md}.

\section{Fine sessione}
\label{sec:03-fine-sessione}

Lorenzo chiude la sessione con un checkpoint esplicito:

\begin{quote}
\texttt{@checkpoint} --- chiudo la sessione.
\end{quote}

L'agente produce:
\begin{itemize}
  \item Il blocco aggiornato per \texttt{\_CONTEXT.md} (Active Task, Token Est.)
  \item Una voce per \texttt{PROGRESS.md} con riepilogo del lavoro
\end{itemize}

Lorenzo copia i due blocchi, fa il commit, chiude il terminale. Domani l'agente saprà esattamente dove erano rimasti.

\section{Il ciclo completo in sintesi}
\label{sec:03-ciclo}

\begin{lstlisting}[language=,caption={Schema del ciclo ADLC per sessione}]
INIZIO SESSIONE
  +- Chiedi conferma contesto → agente legge _CONTEXT.md

DURANTE LA SESSIONE
  +- Task LOW     → esegue + notifica
  +- Task MEDIUM  → propone piano → approvi → esegue
  +- Task HIGH    → HALT → spiega rischi → confermi → esegue
  +- @checkpoint  → ogni 3-5 azioni significative

FINE SESSIONE
  +- @checkpoint → aggiorna _CONTEXT.md + PROGRESS.md → commit
\end{lstlisting}

\section*{\LabelSummary}

\begin{itemize}
  \item Il \textbf{bootstrap} assicura che l'agente abbia letto \texttt{\_CONTEXT.md} prima di qualsiasi lavoro.
  \item I task \textbf{LOW} vengono eseguiti direttamente; i \textbf{MEDIUM} richiedono approvazione del piano; i \textbf{HIGH} richiedono conferma esplicita dopo un HALT.
  \item Il \textbf{checkpoint} ogni 3--5 azioni fissa un punto fermo e mantiene i file di stato aggiornati.
  \item La \textbf{fine sessione} è un checkpoint esplicito che produce gli aggiornamenti da committare.
\end{itemize}

Nei prossimi capitoli esploriamo in dettaglio i meccanismi che rendono possibile tutto questo, a partire dal cuore del sistema: \texttt{\_CONTEXT.md}.
